%%%%%%%%%%%%%%%%%%%%%%%%%%%%%%%%%%%%%%%%%%%%%%%%%%%%%%%%%%%%%%%%%%%%%%%%%%%%%%%
%  Document macros.
%%%%%%%%%%%%%%%%%%%%%%%%%%%%%%%%%%%%%%%%%%%%%%%%%%%%%%%%%%%%%%%%%%%%%%%%%%%%%%%

% --- Project metadata — EDIT THESE, they feed the title page -------------
\newcommand{\thesistitle}{An AI-Powered Lost and Found Platform: Multimodal
  Matching of Item Reports Using Text and Image Embeddings}
\newcommand{\thesisspecialty}{GSI/ISIL}
\newcommand{\thesissupervisor}{TODO --- Supervisor name}
\newcommand{\thesisstudentone}{TODO --- Your full name}
\newcommand{\thesisstudenttwo}{}          % leave empty if working alone
\newcommand{\thesisstudentthree}{}        % leave empty if working alone
\newcommand{\thesisyear}{2025/2026}
\newcommand{\projectname}{LaSauce}

% --- "TO WRITE" boxes ----------------------------------------------------
%  Visible placeholders for the sections that cannot come from the codebase.
%  Flip \showtodosfalse in main.tex to hide every one of them at once.
\newif\ifshowtodos
\newcommand{\towrite}[1]{%
  \ifshowtodos
    \par\medskip\noindent
    \fcolorbox{todoborder}{todobg}{%
      \begin{minipage}{\dimexpr\linewidth-2\fboxsep-2\fboxrule\relax}
        \small\setstretch{1.15}%
        \textbf{\textcolor{todoborder}{TO WRITE.}}\ #1
      \end{minipage}%
    }%
    \par\medskip
  \fi
}

%  Lightweight variant for table cells, where the boxed version's \par and
%  \medskip would fight with the cell's own paragraph handling.
\newcommand{\tofill}[1][to fill in]{%
  \ifshowtodos\textcolor{todoborder}{\small\itshape #1}\fi
}

% --- Figures -------------------------------------------------------------
%  Renders the real image once it exists in figures/, and a labelled grey box
%  until then — so the document always compiles and every gap is visible.
\newcommand{\figplaceholder}[2][\linewidth]{%
  \begingroup
  \setlength{\fboxsep}{0pt}%
  \fcolorbox{rulegray}{black!4}{%
    \begin{minipage}[c][0.24\textheight][c]{\dimexpr#1-2\fboxrule\relax}
      \centering
      {\ttfamily\small figures/#2}\\[0.8em]
      {\itshape\small screenshot not added yet}
    \end{minipage}%
  }%
  \endgroup
}

%  \screenshot[width]{file.png}{Caption}{label}  ->  \ref{fig:label}
\newcommand{\screenshot}[4][0.88\linewidth]{%
  \begin{figure}[htbp]
    \centering
    \IfFileExists{figures/#2}%
      {\includegraphics[width=#1]{figures/#2}}%
      {\figplaceholder[#1]{#2}}%
    \caption{#3}
    \label{fig:#4}
  \end{figure}%
}

% --- Inline shorthands ---------------------------------------------------
%  \code steps down one size to compensate for the larger x-height of the
%  typewriter face in running text. Inside an already-small context (a
%  \scriptsize or \tiny diagram node) that step *up* to \small would overflow
%  the box, so diagrams use \dcode, which inherits the surrounding size.
\newcommand{\code}[1]{\texttt{\small #1}}
\newcommand{\dcode}[1]{\texttt{#1}}
\newcommand{\dq}[1]{\enquote{#1}}
